\section{Interpretabilidad con Grad-CAM}
\label{sec:gradcam}

\textbf{Método (entregable 5).} Grad-CAM~\citep{selvaraju2017gradcam} se calcula así:
\begin{itemize}
  \item se toma el corte de mayor atención del hallazgo, $\arg\max_d a_{d,l}$;
  \item se usan los mapas de activación $A^k$ de \texttt{layer4} en B (7$\times$7) o de la
        última ReLU antes del \emph{pooling} en A (14$\times$14);
  \item se pondera cada mapa por el gradiente medio del \emph{logit} del hallazgo,
        $\mathrm{ReLU}(\sum_k \alpha_k A^k)$;
  \item el mapa se interpola a 224$\times$224 y se superpone al corte.
\end{itemize}
El estudio completo pasa por la red, así que la atención participa en el gradiente. Para
cinco hallazgos mostramos los tres dorados positivos con mayor probabilidad y el falso
positivo más seguro. La Figura~\ref{fig:gradcam} muestra un acierto por hallazgo.

\begin{figure}[t]
  \centering
  \includegraphics[width=\linewidth]{figures/gradcam_B.png}
  \caption{Grad-CAM del modelo B en un dorado positivo por hallazgo, en el corte elegido por
    la atención. Fila superior: corte sagital; fila inferior: mapa superpuesto. El derrame
    (líquido suprapatelar) y la contusión (edema en el cóndilo femoral) se localizan bien. En
    el LCA el mapa cae sobre el derrame y no sobre el ligamento; en el menisco medial, la
    atención elige un corte periférico, fuera de la articulación.}
  \label{fig:gradcam}
\end{figure}

\textbf{Resultados en B.} Revisamos visualmente 20 casos de B, cuatro por hallazgo:
\begin{itemize}
  \item \textbf{Plausible: derrame y contusión.} En los tres aciertos de derrame el calor
        cubre la colección de líquido del receso suprapatelar. En contusión cae dentro del
        cóndilo femoral, sobre el edema óseo moteado.
  \item \textbf{Signos secundarios: LCA.} El calor nunca está en la escotadura intercondílea.
        Los cortes elegidos son laterales y el mapa cae sobre el derrame o el hueso anterior,
        dos signos que acompañan a la rotura: entre los dorados, el 75\,\% de las roturas del
        LCA tienen derrame, frente al 60\,\% global.
  \item \textbf{No plausible: meniscos.} Los mapas no caen sobre el menisco. En el medial,
        los cuatro casos eligen cortes del extremo de la pila (3/24, 3/30, 30/33 y 31/35),
        que muestran grasa y músculo.
\end{itemize}

\textbf{Atención por corte.} En los 58 dorados, la atención para el LCA cae en el tercio
central de la pila solo en el 2\,\% de los casos, aunque el ligamento está en el centro. Los
dos meniscos se atienden en extremos opuestos solo en el 43\,\% de los casos, cerca del azar.

El modelo A, en cambio, concentra la atención en el centro para casi todos los hallazgos
(LCA 67\,\%, fractura 79\,\%), con mapas casi idénticos entre hallazgos del mismo estudio:
aprendió un detector genérico de ``rodilla anormal''.

\textbf{Lectura.} B es anatómicamente coherente en los hallazgos que más resaltan en una
secuencia sensible a líquido, y son también aquellos en que mejor rinde. En estructuras finas
falla: un menisco ocupa menos de una celda del mapa de 7$\times$7. Allí el modelo recurre a
correlaciones. Clínicamente, un detector de rotura del LCA que se apoya en el derrame
fallaría ante una rotura sin derrame, un riesgo que la AUC de 0,71 no deja ver.
